\chapter{Discussion}

\section{Rationale for Studies Performed in the Dissertation}
The studies performed in the dissertation were designed to give insight to the discrete attenuation mechanisms for the \gls{yfv} 17D strain vaccines, which are poorly understood to date, despite widespread and longstanding use of the vaccine exceeding 600 million doses administered \parencite{Monath:2013cu}. Very limited work has been performed on this topic by means of reverse genetics systems as these studies are complicated by a lack of immunocompetent small animal models in which to screen putative attenuating alleles for viscerotropic disease. The quasispecies paradigm offer advantages by which to first explore the selection pressures encountered by the \gls{wt} and vaccine strains of \gls{yfv}, and by doing so narrowly identify a stable genotype under selection in the vaccine.  At the start of this work, viral population structure was completely unexplored for \gls{yfv} in the laboratory setting, and the few datasets in open literature were of insufficient coverage or quality to observe the deep viral populations at any reasonable level of complexity \parencite{Victoria:2010kf}. 
Studies of viral diversity are frequently presented \emph{ex nihilo}, or without proper biological context for the features or structures described therein. Consequently, the central motivating insight for development of the presented studies was that, by using standardized vaccine strains of well-described properties, the association of standard phenotype with \emph{in silico} datasets is accomplished with rigorous, naturally-arising controls.  

Historically, the precedent for a diversity paradigm of \gls{yfv} predates the widespread use of 17D or French vaccine types. Before standardization of 17D and the creation of the seed-lot system, stability of \gls{yfv} vaccine strains was frequently investigated by serial passage studies in animals where the properties of viscerotropism and neurotropism were observed to be partially unlinked, meaning that each phenotype could be selected independently by passage \emph{in vivo} \parencite{FINDLAY:1935vu}. The low viremia and mild infection caused by 17D was later shown to be associated with restriction of the vaccine virus from tissue dissemination, which is a consequence of innate immunity \parencite{Erickson:2013kv}; conversely, productive \gls{wt} \gls{yfv} infection in nonhuman primates results in high viremia and eventual tropism for the liver \parencite{Monath:1981tc}.  If tissue tropism and compartment traversal in some manner is associated with \gls{wt} infection (e.g., viscerotropism and neurotropism phenotypes), the mechanisms underlying this are almost certainly exerted at the subpopulation level, meaning that some alleles are carried in the population may occupy divergent fitness niches within a single host or infection event.

Because many of the vaccine strains available in the study arise from controlled or standardized preparations for which precise information is available, a number of open research questions that derive insight from classical vaccinology of \gls{yfv} or rational understanding of attenuation mechanisms were then rendered amenable to experimentation. The \textbf{overall objective} of the studies was to elucidate viral diversity and population structure for the \gls{yfv} 17D vaccine, and in doing so to provide insight to the mechanism(s) of attenuation of the vaccine. The \textbf{central hypothesis} was that specific features of population structure would be observed to differentiate \gls{wt} \gls{yfv} from the live-attenuated vaccine strain 17D. 

\section{Main Findings and Future Directions}
\subsection{Chapter 3}
In chapter 3, a deep sequencing study was performed to compare the population structures of the \gls{yfv} 17D-204 vaccine with the \gls{wt} parental strain Asibi, from which the vaccine was derived. This study was performed to establish a baseline comparative relationship between the two strains related by precise documentation of serial attenuating passage, under a hypothesis that, besides providing confirmation of the known consensus genotype of the vaccine (which is well-described and used in production standards), sequence subpopulations would be resolved that uniquely classify and differentiate 17D from the parental strain.  By a number of measurements of diversity and by specific sequence content, this hypothesis was supported by the dataset. First, when variant subpopulations in the strains were analyzed, greater numbers of \glspl{snv} were apparent in the \gls{wt} strain, and resolved \glspl{snv} were observed in the parental strain in overall greater frequency than those in the vaccine. Second, a subset of identified variant populations were observed in the Asibi strain bearing overlapping (intersecting) identity with the vaccine whereas the converse was not observed in 17D, indicating that it is possible that some portion of the vaccine genotype is maintained naturally in the \gls{wt} virus, either by random effects or by selection in nature. Third, these patterns were observed by numerous independent measurements of diversity for the two strains. In the light of these findings, it was considered likely that a lack of diversity and attendant lack of subpopulation content would be a typical feature of live-attenuated \gls{yfv} vaccine strains, and in theory other live-attenuated viral vaccines would contain similar features. Some of these questions were addressed in subsequent chapters, especially those of repeatability for subpopulations observed in the vaccines. Selection pressures maintaining the apparent population structure of Asibi in the preparation used (or of other \gls{wt} strains) were not addressed; Asibi was chosen exclusively by virtue of the known and direct relationship between this strain and the vaccine derivatives, in addition to the low passage history for the isolate as maintained in an institutional collection (\gls{wrceva}, \gls{utmb}). However, further the investigation of these features across a wide panel of \gls{wt} \gls{yfv} isolates would be of great use to establish a rigorous baseline for \gls{wt} \gls{yfv} diversity; this would ideally be performed in such a manner as to associate viral populations with particular host ecologies, tissue tropisms, or disease outcomes. The principal context of this study was, however, only limited to the properties of the vaccine.

\subsection{Chapter 4}
In chapter 4, a deep sequencing study was performed to resolve the population structures of a set of phenotypically stable 17D vaccine strains, in this case a set of rare isolates for the \gls{yfv} vaccine 17D. The panel of strains used in the study represents the most complete assembly of genetic relationships for the 17D lineage to date. Lineage and passage relationships for 17D are well described, forming the basis of a seed-lot system that is used to regulate the manufacture of vaccine lots. Although some previous efforts have been undertaken to describe the divergence of 17D vaccines, to this point this has been performed using a combination of consensus sequencing, Bayesian phylogenetics, and classical, \emph{in vitro} methods such as plaque size, or classification of vaccine-specific epitope content \parencite{Barban:2007hi, Stock:2012gp}. 17D is considered highly stable both \emph{in vitro} and \emph{in vivo} methods, and is (with very few benign exceptions between substrains) phenotypically stable with respect to immunogenicity and reactogenicity. Reasonably, it was hypothesized that the previously observed features of stability would be associated with stable population structure of the strains considered. This hypothesis was supported by several findings. First, the diversity profiles of the vaccines were largely similar; statistical divergence of strain diversity was typically observed for strains with distant relationships along the passage lineage, an expected result that is suggestive of some marginal genetic drift of the vaccine; selection analysis by $dN/dS$ revealed the expected pattern of broad negative selection across the vaccine genomes, meaning that marginal differences in subpopulation diversity are most likely attributed to random effects at noncoding sites. Second, the genetic divergence was computed based on the relative frequency of nucleotides in the alignments for each vaccine strain (\gls{rmsd}); however, the use of \gls{rmsd} to generate a fine-scale pairwise divergence measurement did not result in the complete resolution of a typical distance phylogeny for the vaccines, a result that was roughly similar to that produced by consensus sequences. Third, the variant populations of the vaccines and seeds contained marginal levels of \gls{wt} content at low frequencies, which in some cases was observed in repeatable patterns between closely related strains. Fourth, one secondary seed strain (\glspl{ndc} 6676), while included in the comparative analyses, was never used for manufacture of vaccine, as the seed failed a monkey neurovirulence test \parencite{Tauraso:1968vg,Monath:1983uk}. This particular strain contained dissimilar population content in reference to the others, in particularly sister strain 6677, with greater evidence of positive selection; although the consensus sequence of this seed strain was unremarkable, this is likely evidence of mishandling leading to subpopulation divergence of sequence content. In this light of these findings, the presented data represents a reasonable framework by which low-level sequence differences of newly produced 17D vaccine seeds or lots could be assessed against a dataset representing the full lineage of standard strains. Reverse genetics of the particular \glspl{snv} in the 6676 strain populations would be useful to understand the contribution of quasispecies to monkey neurovirulence. Similarly, chimeric Asibi/17D viruses generated by reverse genetics to examine the influence of one or more of the amino acid substitutions that separate Asibi and 17D would provide excellent material to investigate the localized affects of diversity due to specific amino acid substitutions. In particular, these studies would be motivated by a reasonable hypothesis that the rates of nucleotide error afforded by the \gls{yfv} polymerase are a significant contributor to the population dynamics of \gls{yfv}, and so mutagenesis studies on the \glspl{ns}5 protein would likely be fruitful.
\subsection{Chapter 5}
In chapter 5, a deep sequencing study was performed using a panel of live attenuated vaccines that were developed and fielded in parallel with the 17D strain aqueous preparations. These “French neurotropic” strains were developed by the Institut Pasteur (Dakar, Senegal) and widely used in West Africa between 1940 and 1980; they were ultimately discontinued by the Institut Pasteur for reasons of safety. \emph{In silico} methods were used to investigate competing hypotheses on the mechanisms influencing the adverse neurotropic phenotype of the vaccine. Significantly, the strains had been previously assayed for neurovirulence in both mice and monkeys; the strains were unequally neurovirulent in both model systems. Primarily, the results show that for \gls{fnv}, high neurovirulence \emph{in vivo} is associated with a homogenous viral subpopulation structure. Strains with low neurovirulence exhibited complex populations, with considerable \gls{wt} sequence content by consensus and by variant subpopulation. As presented, the findings represent the first analysis of population structure for a set of vaccine strains known to have an unacceptable safety profile. Although some previous efforts had indicated that highly virulent strains would exhibit a greater level of diversity relative to poorly virulent strains, the converse was observed in this study. In the light of these findings, it is likely that the neurotropic adverse reactions caused by \gls{fnv} were the result of the initial adaptation and production of the vaccine in mouse brain, leading to greater propensity of the vaccine strain to enter or productively infect human central nervous tissues. Also most significantly, it is demonstrated that subtle divergence of population structure in \gls{yfv} vaccine preparations may be associated as a continuous quantity with other key measurements of phenotype. Further work will parse the effects of discrete population destabilization for \gls{fnv}, leading to a definitive model of attenuation for \gls{yfv} in nervous tissue.
\subsection{Chapter 6}
In chapter 6, a deep sequencing study was performed to compare the prototype strain Asibi (identical to that used in chapter 3) to that of a strain that was adapted from Asibi by serial passage to productively infect Syrian golden hamsters, causing viscerotropic disease. The adapted strain, called Hamster/P7, was previously reported and as such a consensus sequence comparison of Asibi and Hamster/P7 was already available in open literature, along with histological confirmation of disease outcomes in the model \parencite{McArthur:2003fv}. Evidence of selection for the genotype of Hamster/P7 was demonstrated by reduction in diversity and specific identifying variant populations that differentiate input and output strains. In this light of these findings, it is substantiated that the viscerotropic pathogenicity of \gls{wt} \gls{yfv} follows an “iceberg” paradigm, in which \gls{yfv} and other arbovirus are thought to cause fulminant disease in only a small minority of all infections. In this context, there is support offered for the “trade-off” hypothesis, in which a virus that inhabits a particularly complex ecological niche with multiple hosts is poorly adapted to infect any of the hosts encountered in \gls{wt} transmission. The short [seven] passage series required to produce a pathogenic strain in hamsters substantiates the likelihood that the viscerotropic genotype exists as a naturally occurring subpopulation in the Asibi strain, as opposed to one arising from purely random effects in the host. In further work, the results of this study will be compared to another adapted models arising from the Asibi strain, especially that of an attenuated strain resulting from passage of Asibi in Hela cells, to provide a common reference point for phenotype divergence \parencite{Dunster:1999ij}. 
\subsection{Chapter 7}
In chapter 7, a deep sequencing study was performed using several strains of \gls{jev}, in which a comparison was made between the live-attenuated vaccine strain SA14-14-2 and the \gls{wt} parental strain SA14 from which the vaccine was derived by serial passage in \gls{phk} cells. The intent for this study was to attempt translation of the findings of chapter 3 into another flavivirus vaccine; in doing so the studies would establish that diversity effects are a common feature in live-attenuated, positive-stranded, \gls{rna} virus vaccines. Although the attenuation paths of \gls{yfv} 17D and \gls{jev} SA14-14-2 differ with respect to pathophysiology and tissue tropism, several similarities to \gls{yfv} were apparent in the \gls{jev} dataset. First, SA14-14-2 was of lower diversity than SA14, in the same manner that 17D is low diversity relative to the Asibi strain. Second, the consensus sequences observed for two vaccine strains analyzed were similar, but not identical to, previously reported genotypes for SA14-14-2; it is important to note that there is some dissimilarity between the few strains of both SA14 and SA14-14-2 reported in public databases. Third, variant subpopulations of two SA14-14-2 strains were very similar, despite a difference of once passage in cell culture. Fourth, changes in the population structure in transition from SA14 to the SA14-14-2 vaccine genotype were highly localized; this suggests a similar profile of stability to \gls{yfv} 17D on limited passage. Changes in both diversity and subpopulation frequencies were clustered almost exclusively in the envelope gene of the vaccine; this substantiates previously reported results that used mutagenesis techniques to demonstrate that envelope residues contain the predominant attenuation determinants of SA14-14-2 \parencite{Gromowski:2015et,Yun:2014jr,Sumiyoshi:1995fl}. This differs from the governing paradigm of attenuation for \gls{yfv} 17D, in which the effects are considered to originate with significant mechanistic influence from residues in multiple gene segments.

\section{Critical Methods Analysis for Strategies Used in the Dissertation}

\subsection{RT-PCR and Isolation of Vaccine Nucleic Acids}
To the greatest extent possible, \gls{yfv} strains used in the dissertation studies were not passaged. The intent for this was to preserve the resolvable population structures of vaccine strains without the application of undue selection pressure and consequent alterations to nucleotide frequency. Thus in many cases, sequencing was performed from \gls{rtpcr} amplicons. In such cases, \gls{rtpcr} amplification may indeed introduce mutational biases, either by mutagenesis or by amplification of (assumed by primer design) majority populations defined by the oligonucleotide set. Several controls were used in the \emph{in silico} procedures to limit the effects of potential PCR biases, and are a typical feature of variant calling procedures. First, the techniques and programs used to identify subpopulations themselves used methods designed to limit these potential biases at they would appear in final results to process downstream in the pipeline. Two of these control methods were used in the dissertation. For example in chapter 3, a heuristic variant analysis is presented [Table \ref{Diversity_stats}] for \gls{yfv} strains Asibi and 17D-204, meaning that variant frequencies are retained from the raw percentage of the called nucleotide in the read alignment. In this method, variants are excluded from the final result if biased by greater than a 90 percent divergence in frequency between forward and reverse reads; this procedure offers some protection against both enhancement of particular nucleotides by specific priming and biases in the pool of random hexamers used in the sequencing chemistry itself. Second, it is significant that in chapters for which the working hypothesis was one of population stability between strains, repeatability of variant structure and \gls{wt} content between individual strains was very frequently observed. In chapter 7, \gls{jev} strains separated by a single \emph{in vitro} passage were observed to have roughly identical major coding variant populations, meaning that in some cases, \glspl{snv} are expected to persist under experimental conditions. In chapters 4, 5, 6, and 7, variants are called using a probabilistic criterion in which nucleotide frequencies are resolved on forward and reverse strands separately, and then the frequencies on opposing strains are tested for statistical evidence that the opposing strands  (Fisher’s test, false discovery rate). Further work would ideally involve the deployment of some type of nucleic acid purification to remove the complex egg matrix and later permit direct interrogation of the \gls{vrna} without intervening \gls{rtpcr} amplification.
\subsection{Diversity Indices}
The use of diversity indices is a common method of data exploration in \gls{ngs} studies \parencite{Nishijima:2012ew,Beck:2014cm}. In the described findings, \gls{se} (Shannon’s index) was deployed to visualize and resolve areas of the viral genomes that may be of greater complexity than others, thus warranting closer inspection.  In many of the direct comparisons of strains, the index was useful to estimate raw stability over a passage series, and was especially useful to differentiate strains of divergent phenotype (chapters 3, 5, 6, 7). The significance of a diversity measurement is limited to the extent that the estimate is actually a proxy for variant content, and so the index is strictly meaningful in reference to resolution of the subpopulation genotype contributing the diversity effect at a particular nucleotide position. Typically this value is aggregated across some desired length of the genome in question, or enumerated for a single nucleotide position. In chapter 3, the use of \gls{se} permitted an initial visualization of the Asibi and 17D-204 strain genomes, which revealed that greater number of variant populations were present in the \gls{wt} strain. In chapter 4, diversity was compared between a number of 17D vaccine strains, recovering limited statistical evidence for divergence of the vaccines. It is significant that in chapter 7, in which the \gls{jev} strains were analyzed from \glspl{vrna} directly (unbiased analyses), patterns of low diversity that were hypothesized in reference to the findings in \gls{yfv} datasets were similarly observed. The chapter 7 results are especially significant to question of whether or not 17D and SA14-14-2 are of similar quasispecies structure; the results for the \gls{jev} dataset indicate that diversity effects may be extremely localized, or may involve a very small set of nucleotide sites bearing pathogenically significant nucleotide content.
\subsection{\emph{In Silico} Methods for Resolution of Population Structure}
As mentioned previously, two methods of variant identification were used, with variant calls performed exclusively by probabilistic methods in later chapters. The selected model (V-Phaser v. 2.0) was originally published in a collaboration of developers at the Broad Institute (Cambridge, MA) and Colorado State University \parencite{Yang:2013hk}. Because the initial validation of the program took place using a set of West Nile virus isolates (Family:\emph{Flaviviridae}, genus:\emph{flavivirus}) it was thought that the accuracy of the model would easily translate to a use case with \gls{yfv}, which is identically classified at the genus level. A typical control on the \emph{in silico} procedures would be to perform a plaque-pick followed by Sanger sequencing of the resultant clones, however the usefulness of this procedure would be extremely limited as in the case of the most stable vaccine strains considered in the dissertation (chapter 4), the typical frequencies of rare variants observed in the viral population are relatively low, in such would frustrate attempts to resolve the population clonally. Also, neither heuristic nor probabilistic methods for variant calling as implemented in the studies were designed to resolve linkage between variant populations, meaning the observation of \glspl{snv} as simultaneous arising from the same template strand is not implied by any analysis in the dissertation. Later work would ideally elucidate the groups of \glspl{snv} that may arise from the same template by \emph{in vitro} methods, but models to reconstruct template linkage from \gls{ngs} data are not mature.
\subsection{Selection Analyses}
Selection was estimated by a $dN/dS$ ratio, which was calculated using a custom Python script, referring to the technique of Morelli \parencite{Morelli:2013jl}. For the initial estimation of the possible mutations available, $N$ and $S$ were derived from the consensus sequence of the alignment. This technique was not used in every chapter, but was used in chronologically later efforts, following development of the analysis script. The most significant patterns arising from the use of the $dN/dS$ ratio in the dissertation were foremost that the \gls{yfv} strains analyzed (chapter 3, 5, 6, 7), summary results for the ratio were typically found to not exceed 1.000, meaning that the genome in question is considered to be predominantly influenced by negative (purifying) selection. This pattern was observed at multiple scales in the \gls{yfv} vaccine and parental genomes, across the entire open reading frames of the samples, across discrete genes, and in local patterns resolved by 20-codon sliding window analyses. Purifying selection is a biologically plausible observation, which is frequently observed phylogenetically for \gls{yfv} and other arboviruses \parencite{Schuh:2011ck,Holmes:2003fq}. Some cases were observed in which $dN/dS$ was >1.000, and this was typically localized, as in the case of both \gls{jev} SA14-14-2/P1 and SA14-14-2/P2 (chapter 7) in which $dN/dS$ of \gls{ns}4B was 1.710 and 1.871, respectively. Also in chapter 7, short windows in the envelope genes of the two vaccine strains were of high $dN/dS$, indicating that some level of positive selection influences the vaccine genotype; it is not known as to whether this is represents an essential instability in the vaccine, but passage-stability tests suggest this may not be cause for concern \parencite{Yang:2014io}. One potential problem in the ratio generation is that it relies on the accuracy of the called consensus, and in a naïve manner refers to this sequence as the true, underlying population from which the analyzed strain was selected; the selection of vaccine genotypes from a completely homogenous underlying population is not likely to occur in any \gls{wt} strain handled for attenuating serial passage in the laboratory. It is significant to observe that For later work, it is likely that by constructing a model of population for multiple, related \gls{wt} strains, the selection pressures influencing the population genotypes of derived vaccines would be rendered in greater accuracy, however this has not been attempted in the dissertation lab or by any other group.
\section{Statistical Analyses}
In chapters 3, 4, 5, 6, and 7, diversity indices are compared for divergence nonparametrically, usually by a Kruskal-Wallis/Mann-Whitney test or equivalent. The rationale behind the choice of test is that entropy values are not normally distributed, instead taking a positive-bounded, long-tailed shape similar in appearance to either the gamma or lognormal distributions. This pattern is likely to arise from a complex process, that involves both the inherent noise of the sequencing chemistry combined with the true diversity effects that are being resolved. Other diversity measurements were deployed in the dissertation to provide confirmation to entropy profiles. In chapter 3, Asibi and 17D-204 strains were triply analyzed by entropy, Simpson’s $1-D$ (a more stringent diversity index), and error rate. In chapters 3 and 5, error rates for the strains were compared by quantile-quantile plots, which in some cases revealed divergence of error between strains of differing phenotype. It chapters using multiple independent measurements to confirm diversity patterns, it is significant that techniques were typically found to be in agreement. In further work, it will be appropriate to develop a parametric test that reliably compares diversity profiles between \gls{ngs} datasets, but this is rarely attempted and the statistical tools available from academic sources are largely focused on resolution of variant structure. In further work, it would be highly useful to attempt a continuous regression of diversity profiles against some other continuous variable, e.g. plaque titer, treatment parameter, or lethality endpoint in an animal model. 
\section{Experimental Issues in Flavivirus Vaccinology That May Benefit From A Diversity Framework}
A new and growing body of literature on the topic of viral diversity, while compelling, represents significant challenges to experimental control, namely the presentation of defined viral populations to experimental systems. Studies of this type are typically constructed in such a manner as to assay viral populations when perturbed by some external selection pressure, e.g. pharmaceutical treatment in the host, alternating passage \emph{in vitro}, or comparative rational mutagenesis of an infectious clone strain. Great difficulty arises in the preparation of controlled viral populations for presentation to an experimental system, whether \emph{in vitro} or \emph{in vivo}. A number of findings suggest that \emph{in vivo} population effects would reveal tissue compartmentalization, which is likely for \gls{wt} \gls{yfv} based on the unnatural and widespread dissemination achieved by the virus when innate immunity is ablated \parencite{Erickson:2013kv}. In order to reduce these findings to the fewest number of variables possible, the ideal model system would test an \emph{in vivo} endpoint against a mixed population of a single hypothetical allele. This is sometimes referred to as a “selection” test, in which a single allele with split population frequency is introduced to an experimental system in order to measure selection of allele frequency conferred by that particular system.
\subsection{Reconciling Population-Level Effects WIth Typical Assessments of YFV Vaccine Attenuation}
Current international technical specifications specify that \gls{yfv} vaccine seeds be clinically qualified before use in manufacture; this testing may include comparative immunogenicity and safety studies in human subjects. In addition to the clinical qualification, technical specifications indicate that nonclinical qualification of seeds may take place with a diversity of methods, which presumably would admit the utililty of \gls{ngs} to resolve the sequence or population structure of a production seed in development, in conjunction with the typical use of the standard monkey neurovirulence assay \parencite{WorldHealthOrganization:2010wo}. The findings of these studies indicate a high level of utility for \gls{ngs} analyses in determining the stability of key attenuating features for these vaccines before manufacturers encounter the expense of further qualification. The results of chapters 3, 4, 5, and 7, particularly indicate that with further effort, \gls{ngs} would be useful in the prediction of viral phenotype to ensure highest safety of the vaccine before introduction to living systems.

